\section{Benchmark Construction}
\dnum{[The numeric data throughout the paper will be iteratively updated, hence they are colored and bold]}
\label{sec:benchmark-construction}

SecDevQA is built in five stages: (i)~systematic selection of source
repositories from the GitHub Advisory Database, (ii)~LLM-assisted detection and
extraction of candidate security question--answer (Q\&A) pairs from issue and
discussion threads, (iii)~manual verification and artifact/hard-fact annotation
of each candidate pair, (iv)~normalization of every verified pair into one or
more self-contained, gradeable evaluation items, and (v)~inductive open coding of
the verified pairs toward a taxonomy of developer security questions. The numbers reported below describe the corpus at the
current stage of construction; the benchmark is being actively expanded along
the same protocol, so absolute counts will grow while the procedure remains
fixed.

\subsection{Repository Selection}
\label{sec:repo-selection}

To avoid hand-picking projects in a way that could bias the question distribution, we draw the source repositories from an external, independent maintained catalogue of security activity: the GitHub Advisory Database. We clone the database (snapshot \dnum{2026-06-04}; \dnum{31{,}213} reviewed advisories) and parse every reviewed GHSA entry, extracting the GitHub repositories referenced by each advisory and, for each repository, the number of advisories that carry a documented fix version. This yields \dnum{9{,}102} distinct repositories that have appeared in at least one reviewed advisory. We then apply three pre-specified filters intended to keep repositories that are
both security-active and practically minable:

\begin{enumerate}
  \item \textbf{Security activity:} at least \dnum{10} advisories with a
        documented fix version, so that the project has a substantive, externally
        catalogued security history;
  \item \textbf{Popularity:} at least \dnum{1{,}000} GitHub stars, as a proxy for
        a real user base whose developers ask questions;
  \item \textbf{Maintenance:} a push within one week (as of \dnum{06-06-2026}) and not archived, so that threads reflect current practice.
\end{enumerate}

Applying the advisory-count filter leaves \dnum{440} repositories; adding the
popularity and maintenance filters leaves \dnum{344}. To enforce language and
ecosystem diversity rather than concentrating on whichever ecosystem happens to
have the most advisories, we rank the survivors by fix-bearing advisory count
within each package ecosystem (npm, PyPI, RubyGems, etc.) and retain the top
repositories per ecosystem, producing a diverse candidate pool spanning
\dnum{ten} ecosystems.

From this pool we are mining repositories incrementally. The current corpus
comprises the \dnum{ten} repositories in Table~\ref{tab:repos}, spanning
\dnum{three} languages (Python, JavaScript, Ruby) and a range of
security-sensitive domains:
HTTP clients (\texttt{requests}, \texttt{urllib3}, \texttt{axios}),
cryptography (\texttt{pyca/cryptography}), image processing (\texttt{Pillow}),
web frameworks (\texttt{express}, \texttt{fastapi}, \texttt{rails}),
authentication (\texttt{node-jsonwebtoken}), and payment integration
(\texttt{stripe-node}). The selection deliberately mixes repositories with a
large catalogued advisory history with widely used libraries in domains where
security questions are common but formal advisories are sparser, so that the
corpus is not skewed toward projects with heavyweight advisory processes
(\S\ref{sec:verification} describes the per-repository cap that enforces this).

\subsection{Detection and Extraction of Q\&A Pairs}
\label{sec:extraction}

For each repository we retrieve the most recent \dnum{2{,}000} issue and
discussion threads and flatten them into ordered comment sequences. We bound the
window deliberately: our goal is not to harvest every security Q\&A pair a
project has ever produced, but to collect a corpus large and diverse enough to
support the taxonomy and evaluation in this study. Recent threads also better
reflect current libraries, advisories, and developer practice.

Locating security Q\&A pairs in this material by hand is impractical. A single
active repository accumulates thousands of issues, the large majority of which
are ordinary functional bugs, feature requests, or support questions with no
security content; security discussions are a sparse minority scattered across
that volume, and confirming whether a given thread contains a genuine
question--answer exchange requires reading it end to end. Manually browsing
every issue across \dnum{ten} repositories---tens of thousands of threads---to
find this minority is prohibitively slow. We therefore use an LLM as a
\emph{prefilter} that reads every thread and surfaces likely candidates, turning
an intractable full-text triage into a tractable verification problem over a
small candidate set. A two-stage LLM pipeline identifies candidate Q\&A pairs.

\textbf{Stage~1 (detection).} An LLM reads each thread and decides whether it
contains a genuine developer security question paired with a substantive answer,
emitting a boolean, a short summary, and a confidence score. Because this stage
is a prefilter feeding human verification, we deliberately prompt it to favor
\emph{recall} over precision: the cost of a false positive is a few seconds of
reviewer time at the verification stage, whereas a false negative silently drops
a real pair from the corpus with no opportunity for recovery. We accordingly set
a permissive \dnum{0.70} confidence threshold and instruct the model to surface
borderline threads rather than discard them. The downstream \dnum{39\%} rejection
rate (\S\ref{sec:verification}) is the expected, acceptable price of this
recall-oriented design.

\textbf{Stage~2 (extraction).} For threads passing Stage~1, a second prompt
extracts the specific question and answer comments, the role of the answerer
(maintainer, contributor, commenter, or self-answer), the artifact types the
answer draws on (\S\ref{sec:annotation}), and any externally verifiable
\emph{hard facts} the answer states---CVE/GHSA/OSV identifiers, fixed version
numbers, fix commit SHAs, fix PR references, and advisory URLs.

This pipeline produced \dnum{XXXXX} candidate pairs across the \dnum{ten}
repositories. We intentionally extract candidates broadly at this stage and rely
on manual verification, rather than the LLM, to make the final inclusion
decision.

\subsection{Manual Verification}
\label{sec:verification}

Every candidate pair is reviewed by the first author through a purpose-built
annotation interface that presents the full thread alongside the extracted
question, answer, role, artifacts, and hard facts. A pair is \emph{accepted}
only if it satisfies the inclusion criteria: the question concerns a security
situation in the developer's own project, the answer is substantive and
information-providing, and the question and answer are attributable to distinct,
identifiable comments. Pairs that are not genuinely security-related
(e.g.\ ordinary functional bugs), that are unanswerable without information
absent from the thread, or that are mis-extracted are \emph{rejected}; the
reviewer records a short reason for each rejection.

To prevent the most security-active projects from dominating the corpus, we cap
the number of verified pairs retained per repository at \dnum{25}. This keeps the
corpus balanced across projects, languages, and security domains rather than
skewed toward the few repositories with the densest advisory history. In the
current corpus the cap does not yet bind---every repository falls at or below it
(the largest, \texttt{requests}, contributes \dnum{24} pairs)---but it fixes the
composition policy as mining scales to more and larger repositories.

Of the \dnum{XXXXX} candidates, \dnum{158} (\dnum{60.8\%}) were accepted and
\dnum{102} were rejected---a rejection rate that, given the recall-oriented
prefilter, is expected and underscores why LLM extraction alone is insufficient
and human verification is required. No single repository exceeds \dnum{15.2\%} of
the verified pairs. The accepted pairs are answered predominantly by project
maintainers (\dnum{104} of \dnum{158}), with the remainder from
contributors~(\dnum{26}), other commenters~(\dnum{18}), and the original poster
self-answering~(\dnum{10}). Threads span \dnum{2020--2026} and contain a median
of \dnum{4} comments (mean \dnum{5.7}, max \dnum{41}).

Of the \dnum{158} verified pairs, \dnum{105} (\dnum{66.5\%}) satisfy the
additional \emph{hard-verifiability} requirement: the answer states at least one
externally checkable fact. This hard subset is the part of the benchmark that
supports fact-grounded automated grading (the remaining \dnum{53} pairs are
retained for taxonomy construction so that the taxonomy is not biased toward only
those question types that yield checkable identifiers). Across the hard subset,
the most common fact types are fixed version numbers (\dnum{47} pairs) and
advisory URLs (\dnum{44}), followed by fix PRs (\dnum{26}), CVE identifiers
(\dnum{24}), fix commits (\dnum{21}), and GHSA identifiers (\dnum{14}). Hard
facts are spot-checked against NVD, GHSA, and OSV at construction time, and the
advisory-database snapshot date is recorded with the release to support future
re-verification.

\subsection{Artifact-Type and Hard-Fact Annotation}
\label{sec:annotation}

Each verified pair carries an \texttt{artifacts\_needed} list recording which
artifact types the maintainer's answer drew on, and a structured
\texttt{hard\_facts} record. The artifact-type distribution
(Table~\ref{tab:artifacts}) shows that answering these questions is rarely a
matter of reading source code alone: documentation and project source dominate,
but advisories, pull-request and commit history, dependency manifests, and
external references (RFCs, registries, upstream issues) each appear in a
non-trivial fraction of answers. This spread is what motivates the per-artifact
context conditions in the evaluation (\S\ref{sec:experimental-setup}): different
question categories lean on different artifact types, and the annotation makes
that dependency measurable.

\subsection{Question Normalization}
\label{sec:normalization}

A verified thread is not yet a gradeable evaluation item. The developer's
question is embedded in a multi-comment exchange, leans on earlier comments
(``the error above,'' ``this issue''), and frequently sits in a thread whose
later comments---or even whose title---already disclose the resolution. Handing
such a raw thread to a model under test would leak the answer and conflate reading
comprehension with security reasoning. We therefore normalize each verified pair
into one or more self-contained \emph{(question, answer)} items through a single
LLM pass that is then human-verified. Every decision in this step is governed by
one constraint---the item must stay faithful to what was really asked and
answered---because authenticity is the benchmark's reason for being; the choices
below are each justified against that constraint.

\noindent\textbf{Self-contained, resolution-free questions.}
The question is reconstructed to stand alone: every reporter-supplied specific
(versions, runtime and operating system, configuration, the exact error and
stack-trace text, and all code or proof-of-concept, kept verbatim) is preserved,
and dangling references are resolved against earlier comments. The question must
not, however, state its own resolution---the fixed version, patch commit or pull
request, or advisory identifier introduced \emph{as the answer} is withheld---so
that the item remains a genuine open question rather than a fill-in-the-blank. We
separate the \emph{subject} a reporter cites as their premise (the CVE they are
worried about, the version they run), which legitimately stays, from the
\emph{fix} that resolves it, which must not appear. To make this auditable rather
than a matter of trust in the rewrite, we (i)~retain the verbatim original
reporter comment alongside the normalized question, so a reader can confirm that
no specific was invented or dropped, and (ii)~run a deterministic leak check that
flags any fix-type hard fact---a fixed version, fix pull request or commit, or
advisory URL---present in the question but absent from the pre-answer reporter
context: a fix the reporter stated earlier is legitimate context, whereas a fix
that could only have come from the answer is a leak. Flagged items are corrected
in the human pass.

\noindent\textbf{Thread-grounded answers.}
The gold answer is a faithful distillation of how the maintainer actually
answered, drawn strictly from the thread; the normalizer is forbidden to add
facts, fixes, or reasoning the thread does not contain. When the maintainer's
answer is terse or a pointer---``fixed in \#932, released 9.0.2''---the gold
answer reflects exactly that, labelled honestly as the maintainer's stated
resolution rather than inflated into a fuller explanation we cannot source. We
deliberately do \emph{not} fetch external advisories or patch diffs to enrich the
gold answer at this stage, as doing so would fabricate ground truth the thread
never contained. Instead, the artifacts the answer points to---advisory URLs and
the fix pull requests or commits---are recorded as \texttt{grounding\_sources};
at grading time the referenced pull-request and commit diffs are resolved and
supplied to the judge under the with-context conditions
(\S\ref{sec:experimental-setup}), so that a terse pointer-answer can still
corroborate a richer but correct model response without our having invented the
reference answer.

\noindent\textbf{Knowledge type: an explicit answerability axis.}
The most consequential annotation we attach records whether answering a pair
requires the repository at all. Each pair is labelled \emph{parametric}---a
security or engineering fact any expert could give from general knowledge or from
a public standard the question itself cites, which an LLM could answer with no
repository or document---or \emph{grounded}---resting on this project's specific
behaviour or releases, or on an external source the thread points to. This
distinction is what keeps the evaluation fair. The common pointer-answer ``fixed
in pull request \#932'' is unanswerable without the repository, so scoring a
no-context model on it would measure memorization of one project's pull-request
numbering rather than security capability. Recording knowledge type per pair
turns the no-context condition into a legitimate capability test on the
parametric pairs, makes the grounded pairs the locus of the artifact-context
analysis (\S\ref{sec:experimental-setup}), and turns high no-context accuracy on
grounded pairs into direct evidence of training-data contamination rather than
reasoning. The label is assigned by the normalizer from the answer's semantics and
then hardened by a deterministic rule: if the answer itself states a fix artifact
also recorded in \texttt{hard\_facts}---matched against the answer text, so that a
reporter-cited subject fact does not trigger it---the pair is grounded regardless
of the model's label, because such an artifact is project- or advisory-specific
information a model cannot produce unaided. The rule errs toward \emph{grounded}
by design: mislabelling a repository-internal answer as parametric would wrongly
admit an unanswerable question into the no-context capability test, whereas the
opposite error merely also routes a pair through the richer context conditions.

\noindent\textbf{Conservative decomposition.}
A single thread sometimes bundles distinct information needs---``is my usage
actually affected?'' and ``how do I remediate?''---each answered separately by the
maintainer. We split a thread into multiple items only when the reporter genuinely
raised distinct questions \emph{and} the maintainer's answer addresses more than
one of them; we never invent a sub-question the thread does not answer. This
restraint is the line between reframing a real implicit ask and fabricating
gradeable-looking questions the thread cannot ground---the latter being precisely
the synthetic-question construction our authenticity claim rejects.

\noindent\textbf{Outcome and verification.}
An initial pass over \dnum{50} verified threads yields \dnum{58} gradeable items,
with \dnum{5} threads decomposed into more than one; \dnum{53} items are grounded
and \dnum{5} parametric. The grounded-heavy split is expected: the
hard-verifiability filter (\S\ref{sec:verification}) selects for answers that rest
on a specific fix or advisory, leaving general-knowledge answers a minority. Every
item retains its verbatim original and is marked for human verification, and
\emph{two researchers independently review every normalized pair}. Reviewing the
full set rather than a sample lets us report inter-annotator agreement (Cohen's
$\kappa$) on the knowledge-type label and on the leak and context-completeness
judgments over the entire benchmark, and reconcile every disagreement before a
pair enters the evaluated set.

\subsection{Open Coding Toward a Question Taxonomy}
\label{sec:open-coding}

To characterize \emph{what} developers ask about, we apply inductive open coding
to the verified pairs. The coding unit is the verbatim Q\&A exchange. An LLM
coder (\texttt{gpt-5.4-mini}, accessed through a uniform LiteLLM interface) is
prompted to assign one or two short noun-phrase codes naming the security
concern at issue, given the issue title and body, the Q\&A summary, the full
thread, and the focal question and answer. As with detection, an LLM-assisted
first pass makes coding tractable and consistent at corpus scale while keeping a
human in control of the final code: the first author then reviewed every pair in
the same annotation interface, free to edit or replace any LLM-assigned code, so
the resulting codebook reflects human judgment rather than model output. The
first-author verification note is deliberately withheld from the coder to avoid
anchoring its codes to the acceptance rationale.

This pass produced \dnum{261} code instances over the \dnum{158} pairs
(\dnum{1--2} per pair), with \dnum{235} distinct codes---a fine-grained
vocabulary that reflects how specific developer security questions are
(e.g.\ ``missing subject alternative names,'' ``HTTPS redirect scheme
downgrade,'' ``transitive dependency vulnerability''). The single most frequent
code, \emph{transitive dependency vulnerability}, accounts for \dnum{12}
instances; the long tail of singleton codes confirms that the corpus is not
dominated by a handful of canned questions.

The next step is axial coding: collapsing the open codes into a smaller set of
higher-level categories that will anchor the per-category analysis in
RQ2--RQ3. Table~\ref{tab:axial} presents a \emph{preliminary, author-proposed}
grouping of the \dnum{235} codes into \dnum{eleven} candidate categories,
reported here to convey the shape of the emerging taxonomy. The two largest
groups---HTTP request and response security (SSRF, redirects, headers, cookies)
and TLS and certificate validation---together cover roughly half the corpus,
consistent with the HTTP-client and web-framework projects mined so far. We
stress that this grouping is not final: category names and boundaries will be
fixed through axial coding with an independent second coder, and we will report
inter-rater agreement (Cohen's $\kappa$) on a held-out sample, together with a
saturation analysis as the corpus grows.

\begin{table}[t]
\caption{Repositories in the current SecDevQA corpus. \textbf{Pairs} = Q\&A
  pairs retained after manual verification; \textbf{Hard} = subset whose answer
  carries at least one externally verifiable fact (CVE/GHSA ID, fixed version,
  fix commit/PR, advisory URL). Stars rounded to thousands (\dnum{2026-06-08}).}
\label{tab:repos}
\centering
\footnotesize
\setlength{\tabcolsep}{4pt}
\begin{tabular}{l l r r r}
\toprule
\textbf{Repository} & \textbf{Lang.} & \textbf{Stars} & \textbf{Pairs} & \textbf{Hard} \\
\midrule
psf/requests              & Python & \dnum{54.0k} & \dnum{24} & \dnum{14} \\
axios/axios               & JS     & \dnum{109k} & \dnum{23} & \dnum{19} \\
expressjs/express         & JS     & \dnum{69.1k} & \dnum{23} & \dnum{17} \\
pyca/cryptography         & Python & \dnum{7.6k} & \dnum{21} & \dnum{13} \\
urllib3/urllib3           & Python & \dnum{4.0k} & \dnum{14} & \dnum{6} \\
stripe/stripe-node        & JS     & \dnum{4.4k} & \dnum{13} & \dnum{5} \\
python-pillow/Pillow      & Python & \dnum{13.6k} & \dnum{12} & \dnum{12} \\
rails/rails               & Ruby   & \dnum{58.5k} & \dnum{11} & \dnum{9} \\
auth0/node-jsonwebtoken   & JS     & \dnum{18.2k} & \dnum{9} & \dnum{5} \\
fastapi/fastapi           & Python & \dnum{99.0k} & \dnum{8} & \dnum{5} \\
\midrule
\textbf{Total (\dnum{10} repos)} & & & \textbf{\dnum{158}} & \textbf{\dnum{105}} \\
\bottomrule
\end{tabular}
\end{table}

\begin{table}[t]
\caption{Artifact types the maintainer's answer drew on, across the \dnum{158}
  verified pairs (a pair may draw on several). Annotations are LLM-extracted and
  spot-checked; they drive the context conditions of the evaluation.}
\label{tab:artifacts}
\centering
\footnotesize
\setlength{\tabcolsep}{5pt}
\begin{tabular}{l r r}
\toprule
\textbf{Artifact type} & \textbf{Pairs} & \textbf{\%} \\
\midrule
documentation (docs, README, changelog) & \dnum{126} & \dnum{79.7} \\
code (repository source) & \dnum{92} & \dnum{58.2} \\
advisory (GHSA / OSV / project advisory) & \dnum{35} & \dnum{22.2} \\
pr\_data (PR diff or discussion) & \dnum{23} & \dnum{14.6} \\
issue\_tracker (linked issues) & \dnum{23} & \dnum{14.6} \\
dependency\_manifest & \dnum{19} & \dnum{12.0} \\
external\_reference (RFC, registry, etc.) & \dnum{14} & \dnum{8.9} \\
commit\_history (log, blame, commit) & \dnum{14} & \dnum{8.9} \\
cve\_cwe\_db (NVD/CWE/CVSS) & \dnum{7} & \dnum{4.4} \\
ci\_logs & \dnum{2} & \dnum{1.3} \\
\bottomrule
\end{tabular}
\end{table}

\begin{table*}[t]
\caption{\emph{Preliminary, author-proposed} axial grouping of the open
  codes into candidate higher-level categories. \textbf{Pairs} counts verified
  pairs touching the category (a two-code pair may fall in two categories, so
  the column sums to more than \dnum{158}). This grouping is a working draft to be
  finalized by axial coding with an independent second coder and inter-rater
  agreement; category names and boundaries are not yet fixed.}
\label{tab:axial}
\centering
\footnotesize
\begin{tabular}{l r p{0.62\linewidth}}
\toprule
\textbf{Candidate category} & \textbf{Pairs} & \textbf{Representative open codes} \\
\midrule
HTTP Request and Response Security & \dnum{40} & raw request body, credential leakage on redirect, CSRF token generation, double submit cookie, automatic redirect handling, absolute URL bypass \\
TLS and Certificate Validation & \dnum{38} & SSL certificate verification failure, OpenSSL version mismatch, ASN.1 parsing error, TLS SNI hostname handling, certificate verification failure, stricter SSL verification \\
Dependency and Supply-Chain Security & \dnum{26} & transitive dependency vulnerability, transitive dependency CVE, compromised package versions, dependency version update, prepare script install failure, false positive dependency vulnerability \\
Authentication and Session Integrity & \dnum{17} & webhook signature verification, unsigned JWT handling, question regarding jwt signing, signature verification bypass, token blacklist, JWT logout invalidation \\
Cryptographic Key and Algorithm Use & \dnum{15} & SM2 key algorithm mismatch, RSA key length requirement, cryptographic algorithm default mismatch, post-quantum support roadmap, encrypted RSA key loading, silent hash digest change \\
Input Handling and Denial of Service & \dnum{15} & memory exhaustion DoS, unsafe boundary generation, ReDoS vulnerability, ReDoS protection, regex route matching, array index limit \\
Vulnerability Report Triage & \dnum{14} & false positive report, scanner false positive, false positive vulnerability report, false vulnerability report, vulnerability database error, advisory database update \\
Patch, Release, and Compliance & \dnum{14} & fixed version announcement, SSRF vulnerability backport, safe version guidance, patched version marking, TLS deprecation support, SSRF fix availability \\
Secure Configuration and Defaults & \dnum{13} & deprecated URL parsing, opt-in security setting, missing config types, android network security config, initializer order misconfiguration, missing rate limiting \\
Platform and Version Compatibility & \dnum{10} & RAT on machine, issue tracker reference, local version parsing, Windows Python compatibility, permission denied error, missing attribute regression \\
Concurrency and Data Integrity & \dnum{7} & global state mutation, thread safety data race, shared block cache, deadlocked transaction state, writes outside transaction, thread safety issue \\
\bottomrule
\end{tabular}
\end{table*}

